

%  TIEMPO PRESENTE  %     %  Explicación Teórica, Referencias, Motivación a realizar la Práctica. Teoría/Ecuaciones que describen los fenómenos que se van a estudiar.  %

    \section{Introducción}
La \txcurs{contaminación del agua} es la alteración de la calidad de grandes cuerpos de agua debido a la introducción de agentes nocivos por parte del hombre que comprometen su uso para el consumo humano, la vida acuática o las actividades ecológicas en general, afectando su equilibrio natural \cita{1-2}. Causa la reducción de servicios ecosistémicos como el agua dulce, un recurso escaso que representa únicamente el $2.5\,$\% del agua total de la Tierra. De este porcentaje, tan solo el $0.3\,$\% se encuentra en forma líquida en su superficie \cita{3} y alrededor del $70\,$\% del consumo humano está destinado únicamente a la agricultura \cita{4}. \\\\
Debido a la gran preocupación que genera la escasez del agua potable a nivel mundial y al aumento indiscriminado de contaminantes presentes en ella producto de actividades humanas, el estudio de diversos métodos para su tratamiento ha ido tomando gran importancia en los últimos años. Uno de ellos es la tecnología de \txcurs{plasma}, una técnica de oxidación avanzada que es compatible con el medio ambiente, ya que presenta una temperatura operativa baja, no requiere la adición de productos químicos y tiene una alta eficiencia energética \cita{5-6}. No obstante, para su uso comercial e implementación a gran escala en el sector industrial, existen diversas dificultades técnicas y económicas a superar, por lo que este campo se encuentra todavía en etapa de investigación. \\\\
El plasma es un estado fundamental de la materia que se caracteriza por la presencia de porciones significativas de partículas cargadas en cualquier combinación de iones o electrones, y es la forma más abundante de materia ordinaria en el universo \cita{7}. Se denomina \txcurs{plasma no térmico (NTP}, por sus siglas en inglés\txcurs{)} o \txcurs{plasma frío} a un tipo de plasma con bajo grado de ionización que no se encuentra en equilibrio termodinámico ya que la temperatura de sus electrones es mucho mayor que la temperatura de sus iones o átomos neutros. En tecnologías de plasma, existe un gran número de reactores de NTP de diferentes geometrías para el tratamiento de distintos contaminantes presentes en agua. Mientras que en algunos reactores se realizan descargas eléctricas en electrodos sumergidos en el agua a remediar, en otros la descarga se realiza en el aire, actuando de manera superficial sobre el agua \cita{8-9}. En este tipo de reactores, la mayoría de la energía eléctrica suministrada produce la aceleración de los electrones libres del aire, cuya energía suele ser de $1 \n{eV}$ y cuya temperatura llega a los $10000 \n{K}$ \cita{10}. Los electrones altamente energéticos de estas descargas desencadenan una serie de procesos de ionización y excitación del plasma que permiten la formación de radicales libres reactivos y especies no térmicas, como oxígeno atómico ($\tx{O}$), radicales hidroxilo ($\por\tx{OH}$), ozono ($\tx{O}_3$), peróxidos y radiación ultravioleta. Este proceso ha mostrado rapidez y eficiencia a la hora de degradar y descomponer diversos compuestos orgánicos como pesticidas, farmacéuticos, productos químicos tóxicos, y la destrucción e inactivación de virus y bacterias, sin generar subproductos dañinos para la salud. \\\\
En el laboratorio de Descargas Eléctricas a Presión Atmosférica, se montó un reactor de plasma trielectró-\\ dico para el tratamiento de aguas contaminadas. En el reactor, se genera una \txcurs{descarga de barrera dieléctrica (DBD)} en el aire ambiente mediante la aplicación de alta tensión alterna entre dos electrodos planos fijados a lados opuestos de un material dieléctrico. A una distancia de unos pocos centímetros, se fija a un canal por el que circula el agua a tratar un tercer electrodo conectado a tierra. Al elevar el potencial del sistema de electrodos DBD con un alto voltaje continuo, los microcanales de plasma no térmico generados en la DBD se propagan hacia el tercer electrodo y alcanzan la superficie del agua, produciéndose especies reactivas como las mencionadas anteriormente \cita{11}. Empleando esta configuración, se trabajó en el tratamiento de soluciones acuosas del herbicida 2,4-D obteniéndose muy buenos resultados en la eliminación del contaminante \cita{12}. \\\\
Para mejorar la eficiencia de estos tratamientos, se propone montar un reactor híbrido, donde en el canal de circulación del agua a tratar se colocan piezas de vidrio recubiertas con dióxido de titanio ($\tx{TiO}_2$) como material fotocatalizador. De esta forma, las piezas de titanio son activadas por la radiación ultravioleta emitida por el mismo plasma producido en el reactor. Cuando el $\tx{TiO}_2$ es iluminado con radiación cuya energía es igual o mayor que el band gap del material ($3,2 \n{eV}$), la absorción de un fotón produce un par hueco-electrón que promueve reacciones de oxidación y reducción en la superficie. Se ha observado que el $\tx{TiO}_2$ bajo iluminación UV presenta una alta eficiencia fotocatalítica en la inactivación de microorganismos y en la oxidación de compuestos orgánicos e inorgánicos, siendo los films de dióxido de titanio cristalinos con alto porcentaje de fase anatasa los que han presentado mayor fotoactividad \cita{13}. \\\\
En este trabajo se optimizó el diseño del reactor y se lo caracterizó eléctricamente midiendo tensión, corriente y evaluando la potencia de operación para diferentes configuraciones. Asimismo, se evaluó el funcionamiento del reactor, en términos de eficiencia de eliminación y rendimiento energético mediante el tratamiento de una solución acuosa de azul de metileno. Se presentan también resultados preliminares combinando el tratamiento de plasma con recubrimientos de $\tx{TiO}_2$ ubicados sobre el canal de circulación del agua en la zona de la descarga.